% ============================================================
%  ПЕРЕЛІК УМОВНИХ ПОЗНАЧЕНЬ, СИМВОЛІВ І СКОРОЧЕНЬ
%  Розшифровуємо при першій згадці в тексті; тут -- зведений перелік.
% ============================================================

\chapter*{Перелік умовних позначень, символів і скорочень}
\addcontentsline{toc}{chapter}{Перелік умовних позначень, символів і скорочень}

\begin{longtable}{p{3.2cm} p{11.5cm}}
\textbf{Скорочення} & \textbf{Розшифрування} \\
\midrule
\endhead
AURA Runtime & центральний рівень оцінювання інформаційного ризику та прийняття safety-рішень \\
E2EE & наскрізне шифрування (end-to-end encryption) \\
PAKE & протокол парольно-автентифікованого обміну ключами (password-authenticated key exchange) \\
aPAKE & augmented PAKE; асиметричний PAKE \\
OPAQUE & augmented PAKE-протокол на основі OPRF \\
OPRF & сліпа псевдовипадкова функція (oblivious pseudorandom function) \\
PQ & постквантовий (post-quantum) \\
KEM & механізм інкапсуляції ключа (key encapsulation mechanism) \\
ML-KEM-768 & постквантовий KEM (FIPS 203), рівень безпеки 3 \\
X25519 & протокол Діффі--Геллмана на кривій Curve25519 \\
Ed25519 & схема цифрового підпису EdDSA на Curve25519 \\
Argon2id & функція хешування паролів із підсиленням пам'яті \\
HKDF & функція виведення ключів на основі HMAC \\
HMAC & код автентифікації повідомлень на основі гешу \\
AEAD & автентифіковане шифрування з асоційованими даними \\
MRAE & misuse-resistant authenticated encryption \\
DH & протокол Діффі--Геллмана \\
Double Ratchet & протокол подвійного «храповика» Signal \\
X3DH / PQXDH & протоколи початкового узгодження ключів Signal (класичний / постквантовий) \\
MLS & Messaging Layer Security (RFC 9420) \\
relay & серверно-сліпий канал нагляду (server-blind supervision relay) \\
guardian endpoint & пристрій опікуна у сценарії контрольованого перегляду \\
ACK & підтвердження доставки/встановлення \\
TTL & час життя запису (time to live) \\
FFI & інтерфейс міжмовної взаємодії (foreign function interface) \\
ДСТУ & національний стандарт України \\
\texttt{PASS} / \texttt{FAIL} / & статуси release-gated оцінювання якості \\
\texttt{INSUFFICIENT\_SUPPORT} / \texttt{BLOCKED} & \\
\end{longtable}
